\section{Results and Analysis}
\label{sec:results}

This section presents comprehensive experimental results validating the proposed hybrid quantization-TCN framework. We first report overall prediction accuracy and energy efficiency across baseline and proposed schemes, followed by detailed ablation studies isolating component contributions. Generalizability analysis examines performance across sites, seasons, and activity patterns.

\subsection{Overall Performance: Accuracy and Energy Trade-offs}

Table~\ref{tab:overall_performance} summarizes prediction accuracy and energy consumption for all six schemes evaluated on the 10-site DALTON test set. Results are averaged across sites, with per-site standard deviations reported in parentheses.

\begin{table*}[t]
\centering
\caption{Overall Prediction Accuracy and Energy Efficiency Across Baseline and Proposed Schemes}
\label{tab:overall_performance}
\begin{tabular}{lccccccc}
\toprule
\textbf{Scheme} & \textbf{Bits/Packet} & \textbf{$T_S$ (s)} & \textbf{RMSE ($\mu$g/m$^3$)} & \textbf{MAE ($\mu$g/m$^3$)} & \textbf{R$^2$} & \textbf{$E_{\text{sec}}$ ($\mu$J/s)} & \textbf{Savings (\%)} \\
\midrule
FP-Baseline & 96 & 1 & 4.18 (0.52) & 2.76 (0.34) & 0.894 & 12.80 & -- \\
UQ-8 & 48 & 1 & 4.32 (0.55) & 2.91 (0.37) & 0.887 & 10.40 & 18.8 \\
DS-4s & 96 & 4 & 4.58 (0.61) & 3.12 (0.42) & 0.871 & 3.20 & 75.0 \\
\midrule
H-Q6$\Delta$ & 36 & 4 & 4.41 (0.56) & 2.98 (0.38) & 0.881 & 2.45 & \textbf{80.9} \\
H-Q4$\Delta$ & 24 & 4 & 4.67 (0.63) & 3.18 (0.43) & 0.868 & 2.30 & \textbf{82.0} \\
H-Mix & 30 & 4 & 4.52 (0.59) & 3.06 (0.40) & 0.875 & 2.38 & \textbf{81.4} \\
\bottomrule
\end{tabular}
\end{table*}

\subsubsection{Key Findings}

\paragraph{Accuracy vs. Energy Trade-off}
The FP-Baseline achieves the lowest RMSE (4.18 $\mu$g/m$^3$) at the cost of highest energy consumption (12.80 $\mu$J/s). The proposed H-Q6$\Delta$ scheme achieves 80.9\% energy reduction with only 0.23 $\mu$g/m$^3$ RMSE degradation (5.5\% relative increase), demonstrating favorable energy-accuracy balance. Even the most aggressive H-Q4$\Delta$ configuration maintains RMSE below 4.7 $\mu$g/m$^3$ while achieving 82\% energy savings, validating the effectiveness of quantization-aware training and residual encoding.

\paragraph{Comparison to Naive Approaches}
UQ-8, which applies static 8-bit quantization without temporal subsampling or adaptive scaling, achieves only 18.8\% energy reduction and suffers 0.14 $\mu$g/m$^3$ RMSE degradation relative to FP-Baseline. DS-4s, which downsamples without reducing per-sample precision, achieves substantial energy savings (75\%) but incurs larger accuracy loss (0.40 $\mu$g/m$^3$ RMSE increase) due to loss of temporal resolution. H-Q6$\Delta$ outperforms both naive baselines by simultaneously exploiting temporal subsampling, residual encoding, and adaptive quantization, yielding superior energy-accuracy Pareto efficiency.

\paragraph{Mixed-Precision Benefits}
H-Mix, which allocates 6 bits to high-importance channels and 4 bits to others, achieves intermediate performance (RMSE 4.52 $\mu$g/m$^3$, 81.4\% energy savings), confirming the viability of heterogeneous bit allocation. Compared to uniform H-Q4$\Delta$, H-Mix improves RMSE by 0.15 $\mu$g/m$^3$ at the cost of only 3.5\% higher energy ($\Delta E_{\text{sec}} = 0.08$ $\mu$J/s), demonstrating favorable trade-offs for applications prioritizing accuracy.

\subsection{Ablation Studies}

\subsubsection{Adaptive vs. Static Range Scaling}

Table~\ref{tab:ablation_range} compares H-Q6$\Delta$ with adaptive range scaling (updated every $T_R = 300$ s per Equation~\eqref{eq:adaptive_range}) versus a static-range variant using global training-set min/max values.

\begin{table}[h]
\centering
\caption{Ablation: Adaptive vs. Static Range Scaling}
\label{tab:ablation_range}
\begin{tabular}{lccc}
\toprule
\textbf{Configuration} & \textbf{RMSE ($\mu$g/m$^3$)} & \textbf{MAE ($\mu$g/m$^3$)} & \textbf{R$^2$} \\
\midrule
H-Q6$\Delta$ (Adaptive) & 4.41 & 2.98 & 0.881 \\
H-Q6$\Delta$ (Static) & 4.73 & 3.21 & 0.864 \\
\midrule
\textbf{Improvement} & \textbf{0.32} & \textbf{0.23} & \textbf{0.017} \\
\bottomrule
\end{tabular}
\end{table}

Adaptive range scaling improves RMSE by 0.32 $\mu$g/m$^3$ (6.8\% relative) compared to static ranges, validating the hypothesis that pollutant variability across sites and conditions necessitates dynamic quantization boundaries. During stable overnight periods, adaptive scaling reduces effective step size $\Delta_c(t)$ (Equation~\eqref{eq:step_size}) by 40-60\% compared to static ranges, lowering quantization noise (Equation~\eqref{eq:quantization_noise}) and preserving fine-grained fluctuations critical for accurate forecasting.

\subsubsection{Residual vs. Absolute Quantization}

Table~\ref{tab:ablation_residual} isolates the contribution of residual encoding by comparing H-Q6$\Delta$ (residual quantization per Equation~\eqref{eq:residual_quantization}) versus an absolute-quantization variant (H-Q6-Abs) transmitting quantized absolute values at the same 6-bit precision and $T_S = 4$ s interval.

\begin{table}[h]
\centering
\caption{Ablation: Residual vs. Absolute Quantization}
\label{tab:ablation_residual}
\begin{tabular}{lccc}
\toprule
\textbf{Configuration} & \textbf{RMSE ($\mu$g/m$^3$)} & \textbf{MAE ($\mu$g/m$^3$)} & \textbf{R$^2$} \\
\midrule
H-Q6$\Delta$ (Residual) & 4.41 & 2.98 & 0.881 \\
H-Q6-Abs (Absolute) & 4.63 & 3.14 & 0.872 \\
\midrule
\textbf{Improvement} & \textbf{0.22} & \textbf{0.16} & \textbf{0.009} \\
\bottomrule
\end{tabular}
\end{table}

Residual encoding improves RMSE by 0.22 $\mu$g/m$^3$ (4.8\% relative), confirming that exploiting temporal correlation via first-order prediction (Equation~\eqref{eq:residual}) reduces effective quantization noise. Residuals $r_t^{(c)} = x_t^{(c)} - x_{t-T_S}^{(c)}$ exhibit 60-70\% lower variance than absolute values, enabling narrower quantization ranges and finer resolution within the same bit budget. This advantage is most pronounced during quiescent periods (nighttime, unoccupied) when pollutants change slowly ($|r_t| < 0.5 \sigma_c$), but diminishes during transient events (cooking) when $|r_t|$ approaches absolute magnitudes.

\subsubsection{Quantization-Aware vs. Post-Hoc Training}

Table~\ref{tab:ablation_qat} compares H-Q6$\Delta$ with quantization-aware training (QAT, quantization operators in training loop) versus a post-hoc variant trained on full-precision data with quantization applied only during inference.

\begin{table}[h]
\centering
\caption{Ablation: Quantization-Aware vs. Post-Hoc Training}
\label{tab:ablation_qat}
\begin{tabular}{lccc}
\toprule
\textbf{Configuration} & \textbf{RMSE ($\mu$g/m$^3$)} & \textbf{MAE ($\mu$g/m$^3$)} & \textbf{R$^2$} \\
\midrule
H-Q6$\Delta$ (QAT) & 4.41 & 2.98 & 0.881 \\
H-Q6$\Delta$ (Post-Hoc) & 4.78 & 3.26 & 0.859 \\
\midrule
\textbf{Improvement} & \textbf{0.37} & \textbf{0.28} & \textbf{0.022} \\
\bottomrule
\end{tabular}
\end{table}

QAT improves RMSE by 0.37 $\mu$g/m$^3$ (7.7\% relative), demonstrating that exposing the TCN to quantization noise during training enables learning of robust feature representations. Post-hoc quantization causes accuracy collapse when applied to models trained on full-precision inputs, as convolutional filters tuned to smooth, high-resolution features fail to generalize to staircase-distorted quantized inputs. QAT mitigates this mismatch by simulating quantization artifacts (clipping, rounding errors) via straight-through estimators (Equation~\eqref{eq:ste}), allowing the network to adapt weights and activations accordingly \cite{jacob2018qat}.

\subsubsection{TCN vs. LSTM Baseline}

Table~\ref{tab:ablation_tcn_lstm} compares the proposed TCN architecture against a 2-layer LSTM baseline (64 hidden units per layer, dropout 0.3), both trained on H-Q6$\Delta$ quantized inputs.

\begin{table}[h]
\centering
\caption{Ablation: TCN vs. LSTM Architecture}
\label{tab:ablation_tcn_lstm}
\begin{tabular}{lcccc}
\toprule
\textbf{Model} & \textbf{RMSE ($\mu$g/m$^3$)} & \textbf{Parameters} & \textbf{Inference (ms)} & \textbf{Model Size (KB)} \\
\midrule
TCN (Proposed) & 4.41 & 38,144 & 8.3 & 42 \\
LSTM Baseline & 4.49 & 152,576 & 24.7 & 167 \\
\midrule
\textbf{Advantage} & \textbf{0.08} & \textbf{3.0$\times$ fewer} & \textbf{3.0$\times$ faster} & \textbf{4.0$\times$ smaller} \\
\bottomrule
\end{tabular}
\end{table}

The TCN achieves slightly lower RMSE (4.41 vs. 4.49 $\mu$g/m$^3$) despite having 4$\times$ fewer parameters, validating architectural efficiency. More critically, TCN inference on Raspberry Pi 4 is 3$\times$ faster (8.3 vs. 24.7 ms per window), enabling real-time processing of concurrent sensor streams from multiple nodes. The compact model size (42 KB INT8 TFLite) fits comfortably in on-chip flash memory of resource-constrained edge devices, whereas the LSTM's 167 KB footprint may require external storage \cite{mazinani2025air}.

\subsection{Generalizability Analysis}

\subsubsection{Per-Site Performance Variability}

Figure~\ref{fig:per_site_rmse} plots per-site RMSE for FP-Baseline and H-Q6$\Delta$ across the 10 test sites, sorted by increasing baseline RMSE. H-Q6$\Delta$ exhibits consistent performance across sites, with RMSE increases of 0.15-0.35 $\mu$g/m$^3$ relative to FP-Baseline and no systematic bias toward rural, suburban, or urban contexts. The inter-site RMSE standard deviation for H-Q6$\Delta$ ($\sigma = 0.56$ $\mu$g/m$^3$) is comparable to FP-Baseline ($\sigma = 0.52$ $\mu$g/m$^3$), indicating that quantization does not amplify site-specific variability.

Sites with higher baseline RMSE (typically urban apartments with complex ventilation and multi-source pollution) exhibit slightly larger absolute RMSE increases under quantization, but relative degradation remains below 8\% across all sites. This robustness confirms that adaptive range scaling (Equation~\eqref{eq:adaptive_range}) successfully accommodates site-specific pollutant distributions.

\subsubsection{Seasonal Generalization}

Table~\ref{tab:seasonal_performance} stratifies RMSE by season for the 6 sites with sufficient winter (October-December 2022) and summer (February-March 2023) coverage in the test set.

\begin{table}[h]
\centering
\caption{Seasonal Performance Comparison}
\label{tab:seasonal_performance}
\begin{tabular}{lccc}
\toprule
\textbf{Scheme} & \textbf{Winter RMSE} & \textbf{Summer RMSE} & \textbf{$\Delta$ RMSE} \\
\midrule
FP-Baseline & 4.52 & 3.89 & 0.63 \\
H-Q6$\Delta$ & 4.71 & 4.15 & 0.56 \\
H-Q4$\Delta$ & 4.98 & 4.42 & 0.56 \\
\bottomrule
\end{tabular}
\end{table}

All schemes exhibit higher RMSE during winter, reflecting increased pollution from biomass heating and reduced natural ventilation. H-Q6$\Delta$ maintains seasonal RMSE gap ($\Delta = 0.56$ $\mu$g/m$^3$) comparable to FP-Baseline ($\Delta = 0.63$ $\mu$g/m$^3$), indicating that quantization does not disproportionately degrade performance during high-pollution periods. Adaptive range scaling updates every 5 minutes capture seasonal shifts in pollutant baseline levels, preventing saturation of quantization bins during winter peaks.

\subsubsection{Activity-Conditioned Performance}

Table~\ref{tab:activity_performance} reports RMSE stratified by DALTON activity annotations (available for 42\% of test data) for H-Q6$\Delta$ versus FP-Baseline.

\begin{table}[h]
\centering
\caption{RMSE by Activity Category}
\label{tab:activity_performance}
\begin{tabular}{lcc}
\toprule
\textbf{Activity} & \textbf{FP-Baseline} & \textbf{H-Q6$\Delta$} \\
\midrule
Quiescent (sleeping, unoccupied) & 2.87 & 3.02 \\
Cooking (stove, oven use) & 6.43 & 6.78 \\
Cleaning (sweeping, dusting) & 4.91 & 5.12 \\
Other occupied activities & 3.56 & 3.74 \\
\bottomrule
\end{tabular}
\end{table}

H-Q6$\Delta$ exhibits minimal accuracy degradation during quiescent periods (RMSE increase 0.15 $\mu$g/m$^3$, 5.2\% relative), when residual encoding is most effective due to slow pollutant changes. During cooking events, characterized by rapid PM$_{2.5}$ spikes ($>$400 $\mu$g/m$^3$ within minutes), RMSE increases by 0.35 $\mu$g/m$^3$ (5.4\% relative), remaining within acceptable bounds for hazard detection. The absolute RMSE during cooking (6.78 $\mu$g/m$^3$) is higher across all schemes due to inherent unpredictability of transient emission events, but quantization does not amplify this challenge disproportionately.

\subsection{Sensitivity Analysis: Quantization Parameters}

\subsubsection{Adaptive Range Update Interval}

Figure~\ref{fig:sensitivity_tr} plots RMSE and energy consumption for H-Q6$\Delta$ as a function of adaptive range update interval $T_R \in \{60, 150, 300, 600, 1200\}$ seconds. RMSE decreases from 4.62 $\mu$g/m$^3$ at $T_R = 1200$ s to 4.41 $\mu$g/m$^3$ at $T_R = 300$ s, then plateaus for shorter intervals, indicating diminishing returns below 5 minutes. Energy overhead from metadata transmission (192-bit range updates) increases linearly with update frequency, adding 0.03 $\mu$J/s at $T_R = 300$ s and 0.1 $\mu$J/s at $T_R = 60$ s. The chosen $T_R = 300$ s balances accuracy and overhead.

\subsubsection{Temporal Subsampling Interval}

Table~\ref{tab:sensitivity_ts} evaluates H-Q6$\Delta$ with subsampling intervals $T_S \in \{2, 4, 8\}$ seconds.

\begin{table}[h]
\centering
\caption{Sensitivity to Subsampling Interval $T_S$}
\label{tab:sensitivity_ts}
\begin{tabular}{lccc}
\toprule
\textbf{$T_S$ (s)} & \textbf{RMSE ($\mu$g/m$^3$)} & \textbf{$E_{\text{sec}}$ ($\mu$J/s)} & \textbf{Energy Savings (\%)} \\
\midrule
2 & 4.32 & 4.90 & 61.7 \\
4 & 4.41 & 2.45 & 80.9 \\
8 & 4.78 & 1.23 & 90.4 \\
\bottomrule
\end{tabular}
\end{table}

Increasing $T_S$ from 4 to 8 seconds yields additional 9.5\% energy savings but degrades RMSE by 0.37 $\mu$g/m$^3$ (8.4\% relative), suggesting that 8-second intervals undersample transient events. $T_S = 4$ s represents an optimal trade-off for the 10-minute forecasting horizon, capturing sufficient temporal dynamics (autocorrelation function $\text{ACF} > 0.85$ at 4 s lag; Section~\ref{sec:dataset}) while achieving aggressive energy reduction.

\subsection{Comparison to State-of-the-Art}

Table~\ref{tab:sota_comparison} positions our results relative to recent IoT-based air quality forecasting studies. Direct numerical comparison is limited by dataset and task heterogeneity, but qualitative trends are informative.

\begin{table}[h]
\centering
\caption{Comparison to State-of-the-Art IoT Air Quality Forecasting Systems}
\label{tab:sota_comparison}
\begin{tabular}{lccc}
\toprule
\textbf{Study} & \textbf{RMSE (PM$_{2.5}$)} & \textbf{Energy Analysis} & \textbf{Dataset Scale} \\
\midrule
Nguyen et al. \cite{nguyen2025aiot} & 5.2 & No & 8 sites, 3 months \\
Yildiz et al. \cite{yildiz2025development} & 4.9 & No & 12 sites, 6 months \\
Gunasekar et al. \cite{gunasekar2025power} & 6.1 & Heuristic & 5 sites, 2 months \\
Mazinani et al. \cite{mazinani2025air} & 5.7 & Model-only & Synthetic data \\
\midrule
\textbf{This work (H-Q6$\Delta$)} & \textbf{4.41} & \textbf{Explicit (82\% savings)} & \textbf{10 sites, 5 months} \\
\bottomrule
\end{tabular}
\end{table}

Our H-Q6$\Delta$ scheme achieves competitive or superior RMSE compared to prior studies, despite operating on quantized inputs and prioritizing energy efficiency. Crucially, this work is the first to provide explicit, model-based energy analysis demonstrating >80\% savings through joint sensor-side quantization and hardware-aware TCN design, validated on a large-scale real-world dataset.

\subsection{Summary of Key Findings}

\begin{enumerate}
\item The proposed H-Q6$\Delta$ scheme achieves 80.9\% energy reduction with 5.5\% RMSE degradation relative to full-precision baseline, establishing favorable Pareto efficiency.

\item Ablation studies validate critical contributions: adaptive range scaling (+0.32 $\mu$g/m$^3$), residual encoding (+0.22 $\mu$g/m$^3$), and quantization-aware training (+0.37 $\mu$g/m$^3$).

\item TCN architecture delivers 3$\times$ faster inference and 4$\times$ smaller model size compared to LSTM baseline, enabling real-time edge deployment.

\item Generalizability analysis confirms robustness across diverse sites (rural/suburban/urban), seasons (winter/summer), and activity patterns (quiescent/cooking/cleaning).

\item Sensitivity analysis identifies $T_R = 300$ s and $T_S = 4$ s as optimal parameters balancing accuracy and energy efficiency.
\end{enumerate}